\chapter{Dokazi}
\label{chap:dokazi}

V srednji šoli se dijaki pri matematiki učijo, \emph{kako} se kaj
izračuna. Na univerzi imajo študentje matematike pred seboj
zahtevnejšo nalogo: poleg \emph{kako} morajo vedeti tudi \emph{zakaj}.
Od njih se pričakuje, da bodo računske postopke znali tudi utemeljiti,
ne pa samo slediti pravilom, ki jih je predpisal učitelj. Razumeti
morajo dokaze znamenitih izrekov in sami poiskati dokaze preprostih
izjav. Da bi se lažje spopadli s temi novimi nalogami, bomo prvi del
predmeta Logika in množice posvetili matematični infrastrukturi:
izjavam, pra\-vi\-lom sklepanja in dokazom. Učili se bomo, kako pišemo
dokaze, kako jih analiziramo in kako jih sami poiščemo.

\section{Kaj je matematični dokaz?}
\label{sec:kaj-je-dokaz}

Osrednji pojem matematične aktivnosti je \emph{dokaz}. Namen dokaza je
s pomočjo točno določenih in vnaprej dogovorjenih \emph{pravil
  sklepanja} utemeljiti neko matematično \emph{izjavo}. Načeloma mora
dokaz vsebovati vse podrobnosti in natanko opisati posamezne korake
sklepanja, ki vodijo do želene matematične izjave. Ker bi bili taki
dokazi zelo dolgi in bi vsebovali nezanimive podrobnosti, matematiki
običajno predstavijo samo oris ali glavno zamisel dokaza. Izkušenemu
matematiku to zadošča, saj zna oris sam dopolniti do pravega dokaza.
Matematik začetnik seveda potrebuje več podrobnosti. Poglejmo si
primer.

\begin{izrek}
  \label{izr:n3-n-deljivo-3}
  Za vsako celo število $n$ je $n^3 - n$ deljivo s~$3$.
\end{izrek}

\noindent
Izkušeni matematik je ta izrek že videl, še sam ne ve kolikokrat,
zato bi le pripomnil, da je število $n^3 - n$ zmnožek treh zaporednih
celih števil.
%
Tak namig bi starejšemu študentu matematike zadoščal, da bi sestavil dokaz:

\begin{quote}
  \begin{proof}
    Ker je $n^3 - n = (n-1) \cdot n \cdot (n+1)$, je $n^3 - n$ zmnožek
    treh zaporednih celih števil, od katerih je eno deljivo s~$3$,
    torej je tudi $n^3 - n$ deljivo s~$3$.
  \end{proof}
\end{quote}

\noindent
Čeprav bi bila
večina matematikov s tem dokazom zadovoljna, bi morali za popoln dokaz
preveriti še nekaj podrobnosti:
%
\begin{itemize}
\item Ali velja $n^3 - n = (n-1) \cdot n \cdot (n+1)$?
\item Ali je izmed treh zaporednih celih števil eno deljivo s~$3$?
\item Ali je zmnožek treh števil deljiv s~$3$, če je eno od njih deljivo s~$3$?
\end{itemize}
%
S srednješolskim znanjem algebre ugotovimo, da je odgovor na prvo
vprašanje pritrdilen. Tudi odgovora na drugo in tretje vprašanje sta
očitno pritrdilna, mar ne? To pa ne pomeni, da ju ni treba dokazati.
Nasprotno, zgodovina matematike nas uči, da moramo prav ">očitne"<
izjave še posebej skrbno preveriti.
%
V idealnem primeru bi bil dokaz zapisan tako, da bi lahko njegovo pravilnost preverili mehansko, z
računalnikom.\footnote{Kogar to zanima, se lahko nauči uporabljati kak \href{https://ncatlab.org/nlab/show/proof+assistant}{dokazovalni pomočnik}, v zadnjem času hitro napreduje \href{https://lean-lang.org}{Lean}.}

Če dokaz izreka~\ref{izr:n3-n-deljivo-3} zapišemo preveč
podrobno, postane dolgočasen in ne\-ra\-zumljiv:

\begin{quote}
  \begin{proof}
    Naj bo $n$ poljubno celo število. Tedaj velja
    %
    \begin{align*}
      n^3 - n
      &= n \cdot n^2 - n \cdot 1 \\
      &= n \cdot (n^2 - 1) \\
      &= n \cdot ((n + 1) \cdot (n - 1)) \\
      &= n \cdot ((n - 1) \cdot (n + 1)) \\
      &= (n \cdot (n - 1)) \cdot (n + 1) \\
      &= (n - 1) \cdot n \cdot (n + 1).
    \end{align*}
    %
    Vidimo, da je $n^3 - n$ zmnožek treh zaporednih celih števil.
    Dokažimo, da je eno od njih deljivo s~$3$. Število $n$ lahko
    enolično zapišemo kot $n = 3 k + r$, kjer je $k$ celo število
    in $r = 0$, $r = 1$ ali $r = 2$. Obravnavajmo tri primere:
    %
    \begin{itemize}
    \item če je $r = 0$, je $n = 3 k$, zato je $n$ deljiv s~$3$,
    \item če je $r = 1$, je $n - 1 = (3 k + 1) - 1 = 3 k + (1 - 1) = 3
      k + 0 = 3 k$, zato je $n-1$ deljiv s~$3$,
    \item če je $r = 2$, je $n + 1 = (3 k + 2) + 1 = 3 k + (2 + 1) = 3
      k + 3 = 3 k + 3 \cdot 1 = 3 (k +1)$, zato je $n+1$ deljiv s~$3$.
    \end{itemize}
    %
    Vemo torej, da je $n - 1$, $n$ ali $n + 1$ deljiv s~$3$.
    Obravnavamo tri primere:
    %
    \begin{itemize}
    \item Če je $n - 1$ deljiv s~$3$, tedaj obstaja celo število
      $k$, da je $n - 1 = 3 k$. V tem primeru je $(n - 1) n (n + 1) =
      (3 k) n (n + 1) = 3 (k n (n + 1))$, zato je $(n - 1) n (n + 1)$
      deljivo s~$3$.
    \item Če je $n$ deljiv s~$3$, tedaj obstaja celo število $k$,
      da je $n = 3 k$. V tem primeru je $(n - 1) n (n + 1) = (n - 1)
      (3 k) (n + 1) = (3 k) (n - 1) (n + 1) = 3 (k (n - 1) (n +
      1))$, zato je $(n - 1) n (n + 1)$ deljivo s~$3$.
    \item Če je $n + 1$ deljiv s~$3$, tedaj obstaja celo število
      $k$, da je $n + 1 = 3 k$. V tem primeru je $(n - 1) n (n + 1) =
      (n - 1) n (3 k) = (n - 1) (3 k) n = (3 k) (n - 1) n = 3 (k (n -
      1) n)$, zato je $(n - 1) n (n + 1)$ deljivo s~$3$.
    \end{itemize}
    %
    V vsakem primeru je $(n - 1) n (n + 1)$ deljivo s~$3$. Ker smo
    dokazali, da je $n^3 - n = (n - 1) n (n + 1)$, je tudi $n^3 - n$
    deljivo s~$3$.
  \end{proof}
\end{quote}

Tovrstni dokazi so res le sami sebi namen.
%
Dober dokaz poleg gole resnice bralcu na razumljiv način posreduje tudi matematične ideje, na katerih sloni dokaz.
%
Ta socialna komponenta dokaza je izrazito odvisna ne le od pisca, ampak tudi od bralca, ki mu je dokaz namenjen.
%


%%%%%%%%%%%%%%%%%%%%%%%%%%%%%%%%%%%%%%%%%%%%%%%%%%%%%%%%%%%%%%%%%%%%%%

\section{Struktura dokaza}

Matematiki v sklopu svojih aktivnosti \emph{konstruiramo} matematične objekte:
%
\begin{itemize}
\item v geometriji konstruiramo trikotnike z ravnilom in šestilom,
\item računamo števke števila $\pi$, s čimer konstruiramo njegove približke,
\item ko rešujemo enačbo, konstruiramo število, ki ji zadošča,
\item konstruiramo elemente množice, na primer $(2, \inl(3)) \in \NN \times (\ZZ + \ZZ)$.
\end{itemize}
%
Tudi na dokaze lahko gledamo kot na konstrukcije, ki jih gradimo s pomočjo \emph{pravil sklepanja}.
%
Le-ta povedo, kakšna so ">pravila igre"< v dokazovanju, ne pa tudi, kako dobro igrati igro dokazovanja.
%
Pri tem predmetu bomo spoznali pravila in podali nekaj namigov, kako konstruirati osnovne dokaze, pri ostalih predmetih pa boste videli ">žive dokaze"<, ki imajo manj tehničnih podrobnosti in se bolj
osredotočajo na komuniciranje idej.

Dokaz ima ugnezdeno strukturo: sestoji iz pod-dokazov, ki sestoje iz nadaljnjih pod-dokazov itn., ki se zaključijo z osnovnimi koraki. Ta struktura iz običajnega besedila ni razvidna, a je prisotna. Za zgled vzemimo naslednji izrek.

\begin{izrek}
  \label{izrek:pqr}%
  Za vse izjave $p$, $q$, $r$ velja
  $(p \lor q) \land r \lthen (p \land r) \lor (q \land r)$.
\end{izrek}

\begin{proof}
  Predpostavimo $p \lor q$ in $r$. Če velja $p$, potem sledi $p \land r$ ter od tod $(p \land r) \lor (q \land r)$. Če pa velja $q$, sledi $q \land r$ ter spet $(p \land r) \lor (q \land r)$.
\end{proof}

\noindent
Zgornji dokaz z dodanimi podrobnostmi in označeno ugnezdeno strukturo je prikazan na sliki~\ref{fig:izrek-pqr}.
%
\begin{figure}[t]
  \centering
  \begin{boxedProof}[numbered]
  \begin{boxFrame}
    \boxLine{Naj bodo $p$, $q$ in $r$ izjave.\label{vr:spremenljivke}}
    \boxLine{Dokažimo $(p \lor q) \land r \lthen (p \land r) \lor (q \land r)$:}
    \boxLine[1]{Predpostavimo $(p \lor q) \land r$.\label{vr:predpostavka}}
    \boxLine[1]{Dokažimo $(p \land r) \lor (q \land r)$:}
    \boxLine[2]{Zaradi~\boxRef{vr:predpostavka} velja $p \lor q$.\label{vr:disjunkcija}}
    \boxLine[2]{Zaradi~\boxRef{vr:predpostavka} velja $r$.\label{vr:r}}
    \boxLine[2]{Uporabimo~\boxRef{vr:disjunkcija} in obravnavamo dva primera:\label{vr:primera}}
    \begin{boxNested}[2]
      \boxLine{Predpostavimo $p$.\label{vr:p}}
      \boxLine{Dokažimo $(p \land r) \lor (q \land r)$:}
      \boxLine[1]{Dokažimo levi disjunkt $p \land r$:}
      \boxItem[2]{$p$ velja zaradi~\boxRef{vr:p},\label{vr:p-levi}}
      \boxItem[2]{$r$ velja zaradi~\boxRef{vr:r}.\label{vr:r-levi}}
    \end{boxNested}
    \begin{boxNested}[2]
      \boxLine{Predpostavimo $q$.\label{vr:q}}
      \boxLine{Dokažimo $(p \land r) \lor (q \land r)$:}
      \boxLine[1]{Dokažimo desni disjunkt $q \land r$:}
      \boxItem[2]{$q$ velja zaradi~\boxRef{vr:q},\label{vr:q-desni}}
      \boxItem[2]{$r$ velja zaradi~\boxRef{vr:r}.\label{vr:r-desni}}
    \end{boxNested}
  \end{boxFrame}
  \end{boxedProof}
  \caption{Dokaz izreka~\ref{izrek:pqr}}
  \label{fig:izrek-pqr}
\end{figure}
%
Zamik vrstice pove, kako globoko v dokazu stojimo: vrstici, ki postavi cilj, sledi njen dokaz eno stopnjo globlje. S škatlo obdamo tisti del dokaza, pri katerem iz zamikov samih ne bi bilo razvidno, kje se konča. Zunanja škatla najprej vpelje spremenljivke $p$, $q$ in $r$, nato pove, kaj bomo dokazali,
temu pa sledi dokaz. Vrstice so oštevilčene, da se lahko nanje sklicujemo. Notranji škatli sta obravnavi dveh primerov.

Ko pišemo ali beremo dokaz, moramo v vsakem trenutku poznati
%
\begin{itemize}
\item \textbf{cilj}: kaj trenutno dokazujemo in
\item \textbf{kontekst}: katere spremenljivke, predpostavke in že dokazana dejstva imamo na voljo.
\end{itemize}
%
Na primer, v vrstici~\boxRef{vr:primera} iz dokaza na sliki~\ref{fig:izrek-pqr} je cilj $(p \land r) \lor (q \land r)$, veljavne so spremenljivke $p$, $q$ in $r$, na voljo imamo predpostavko~\boxRef{vr:predpostavka}, dokazali pa smo že dejstvi~\boxRef{vr:disjunkcija} in~\boxRef{vr:r}.
%
Ko napravimo korak v dokazu, mora biti utemeljen z enim od pravil sklepanja. Dokaz je
popoln, ko smo utemeljili vse pod-dokaze.
%
V nadaljevanju spoznajmo pravila sklepanja in druge korake, ki tvorijo dokaz.

\paragraph*{Vpeljava spremenljivke}
%
Nekatera pravila sklepanja zahtevajo, da v dokazu \df{vpeljemo spremenljivko}. To storimo tako, da jo naznanimo skupaj z njeno domeno:
%
\begin{quote}
  \nls{Naj bo $x \in A$.}
\end{quote}
%
S tem postane $x$ veljavna spremenljivka, ki jo smemo uporabljati.
%
O njej vemo le to, da je element množice $A$, ne vemo pa, kateri element -- pravimo, da je $x$ \df{prosta spremenljivka}. V matematičnih besedilih boste zasledili tudi naslednje fraze, ki uvedejo prosto spremenljivko:
%
\begin{itemize}
\item \nls{Naj bo $x \in A$ poljuben.}
\item \nls{Fiksirajmo $x \in A$.}
\item \nls{Izberimo $x \in A$.}
\end{itemize}
%
\emph{Pozor:} besedi ">fiksirajmo"< in ">izberimo"< bi dali komu misliti, da si lahko izbere točno določen element po lastni volji, a si je bolje predstavljati, da je~$x$ ">neznan"< ali ">nedoločen"<.

Spremenljivka ima območje veljavnosti, ki je omejeno s pod-dokazom, v katerem jo vpeljemo.

\begin{zgled}
  V dokazu na sliki~\ref{fig:izrek-pqr} vpeljemo spremenljivke $p$, $q$ in $r$ v vrstici~\boxRef{vr:spremenljivke}. Veljavne so v celotnem dokazu.
\end{zgled}

\emph{Pozor:} spremenljivk ne vpeljemo kar tako, ampak le takrat, ko uporabimo pravilo sklepanja, ki to zahteva.


\paragraph*{Vpeljava predpostavke}
\label{sec:vpelj-predp}

Nekatera pravila sklepanja zahtevajo, da \df{vpeljemo predpostavko}~$\phi$ (ali \df{hipotezo}), se pravi, da se pretvarjamo, da izjava~$\phi$ velja. V takem primeru zapišemo:
%
\begin{quote}
  \nls{Predpostavimo $\phi$.}
\end{quote}
%
Včasih vpeljemo predpostavko tudi s \nls{če velja $\phi$:}.

Predpostavka ima svoje območje veljavnosti, ki je omejeno s pod-dokazom, v katerem jo vpeljemo.

\begin{zgled}
  V dokazu na sliki~\ref{fig:izrek-pqr} v vrstici~\boxRef{vr:predpostavka} predpostavimo $(p \lor q) \land r$, območje veljavnosti je do konca dokaza. V vrstici~\boxRef{vr:p} predpostavimo $p$, območje veljavnosti pa je prva notranja škatla. V vrstici~\boxRef{vr:q} predpostavimo~$q$, območje veljavnosti pa je druga notranja škatla.
\end{zgled}

Kadar se skličemo na že znano izjavo ali na predpostavko, pravimo, da smo jo \df{uporabili}.

\begin{zgled}
  V dokazu na sliki~\ref{fig:izrek-pqr} smo predpostavko~\boxRef{vr:predpostavka} uporabili v vrsticah~\boxRef{vr:disjunkcija} in~\boxRef{vr:r}.
  V vrsticah~\boxRef{vr:p-levi} in~\boxRef{vr:r-levi} smo uporabili predpostavko~\boxRef{vr:p} in dejstvo~\boxRef{vr:r}.
  V vrsticah~\boxRef{vr:q-desni} in~\boxRef{vr:r-desni} smo uporabili predpostavko~\boxRef{vr:q} in dejstvo~\boxRef{vr:r}.
\end{zgled}

\emph{Pozor:} predpostavk ne vpeljemo kar tako, ampak le takrat, ko uporabimo pravilo sklepanja, ki to zahteva.

\paragraph*{Okrajšava}
\label{sec:okraj}

V dokazu smemo kadarkoli \df{vpeljati okrajšavo}, na primer:
%
\begin{quote}
  \nls{Naj bo $\varphi \defeq \frac{1 + \sqrt{5}}{2}$.}
\end{quote}
%
V nadaljevanju dokaza smemo uporabljati~$\varphi$ in dejstvo, da je enak $(1 + \sqrt{5})/2$.

\begin{vaja}
  Ali je v zvezi z $(1 + \sqrt{5})/2$ kaj posebnega?
\end{vaja}

\paragraph*{Izpeljava dejstva}
\label{sec:izpelj-dejstva}

V dokazu smemo kadarkoli \df{izpeljati dejstvo}, kar pomeni, da utemeljimo neko izjavo, ki jo nameravamo kasneje uporabiti. Če je utemeljitev dejstva~$\phi$ kratka, zapišemo:
%
\begin{quote}
  \nls{Zaradi \dots\ velja $\phi$.}
\end{quote}
%
Če utemeljitev zahteva celoten pod-dokaz, ga podamo:
%
\begin{center}
  \begin{boxedProof}
    \begin{boxFrame}
      \boxLine{Velja $\phi$:}
      \boxLine[1]{\brac{dokaz $\phi$}}
    \end{boxFrame}
  \end{boxedProof}
\end{center}

\begin{zgled}
  V dokazu na sliki~\ref{fig:izrek-pqr} nismo vpeljali nobene okrajšave, smo pa v vrsticah~\boxRef{vr:disjunkcija}
  in~\boxRef{vr:r} izpeljali dejstvi $p \lor q$ in~$r$. Dejstvo~\boxRef{vr:disjunkcija} smo uporabili v vrstici~\boxRef{vr:primera}, dejstvo~\boxRef{vr:r} pa v vrsticah~\boxRef{vr:r-levi} in~\boxRef{vr:r-desni}.
\end{zgled}

Posvetimo se še pravilom sklepanja, ki jih delimo na:
%
\begin{itemize}
\item pravila sklepanja za logične operacije:
  \begin{itemize}[nosep]
  \item pravila vpeljave, s katerimi dokazujemo izjave,
  \item pravila uporabe, s katerimi uporabljamo predpostavke in dejstva;
  \end{itemize}
\item pravila sklepanja za enakost;
\item pravilo o izključenem tretjem in dokaz s protislovjem.
\end{itemize}

\section{Pravila vpeljave in uporabe}

Vsaka logična operacija ima svoja pravila vpeljave in uporabe, s katerimi dokazujemo še nedokazane izjave in uporabljamo že dokazane izjave.

\paragraph*{Konjunkcija}
\label{sec:konjunkcija-pravila}

Konjunkcijo $\phi \land \psi$ dokažemo tako, da posebej dokažemo~$\phi$ in posebej~$\psi$.
Pravilo vpeljave torej predstavimo takole:
%
\begin{center}
  \begin{boxedProof}
    \begin{boxFrame}
      \boxLine{Dokažimo $\phi \land \psi$:}
      \begin{boxNested}
        \boxLine{Dokažimo $\phi$:}
        \boxLine[1]{\brac{dokaz $\phi$}}
      \end{boxNested}
      \begin{boxNested}
        \boxLine{Dokažimo $\psi$:}
        \boxLine[1]{\brac{dokaz $\psi$}}
      \end{boxNested}
    \end{boxFrame}
  \end{boxedProof}
\end{center}
% 
Kaj pa, če že vemo, da velja $\phi \land \psi$? Tedaj lahko izpeljemo~$\phi$ in izpeljemo~$\psi$, zato imamo dve pravili uporabe:
%
\begin{center}
\begin{boxedProof}[reserved]
\begin{boxFrame}
  \boxLineTag{$k$}{velja $\phi \land \psi$}
  \boxGap
  \boxLine{zaradi~($k$) velja $\phi$}
\end{boxFrame}
\qquad\qquad
\end{boxedProof}
%
\begin{boxedProof}[reserved]
\begin{boxFrame}
  \boxLineTag{$k$}{velja $\phi \land \psi$}
  \boxGap
  \boxLine{zaradi~($k$) velja $\psi$}
\end{boxFrame}
\end{boxedProof}
\end{center}

\paragraph*{Implikacija}
\label{sec:implikacija-pravila}

Implikacijo $\phi \lthen \psi$ dokažemo tako, da predpostavimo~$\phi$ in dokažemo~$\psi$.
Pravilo vpeljave torej predstavimo takole:
%
\begin{center}
  \begin{boxedProof}
    \begin{boxFrame}
      \boxLine{Dokažimo $\phi \lthen \psi$:}
      \boxLine[1]{Predpostavimo $\phi$.}
      \boxLine[1]{Dokažimo $\psi$:}
      \boxLine[2]{\brac{dokaz $\psi$}}
    \end{boxFrame}
  \end{boxedProof}
\end{center}
%
Če vemo, da velja $\phi \lthen \psi$, in vemo, da velja~$\phi$, lahko izpeljemo~$\psi$.
%
Pravilo uporabe za implikacijo, ki se imenuje tudi \df{modus ponens}, se torej glasi:
%
\begin{center}
  \begin{boxedProof}[reserved]
    \begin{boxFrame}
      \boxLineTag{$k$}{velja $\phi \lthen \psi$}
      \boxGap
      \boxLineTag{$m$}{velja $\phi$}
      \boxGap
      \boxLine{zaradi~($k$) in~($m$) velja $\psi$}
    \end{boxFrame}
  \end{boxedProof}
\end{center}

\begin{zgled}
  Naj bosta $\phi$ in $\psi$ izjavi. Takole dokažemo, da iz $\phi \land \psi$ sledi $\psi \land \phi$:
  %
  \begin{center}
  \begin{boxedProof}[numbered]
    \begin{boxFrame}
      \boxLine{Dokažimo $\phi \land \psi \lthen \psi \land \phi$:}
      \boxLine[1]{Predpostavimo $\phi \land \psi$.} \label{box:conj-comm-1}
      \boxLine[1]{Dokažimo $\psi \land \phi$:} \label{box:conj-comm-2}
      \begin{boxNested}[1]
        \boxLine{Dokažimo $\psi$:}
        \boxLine[1]{Zaradi \boxRef{box:conj-comm-1} velja $\psi$.} \label{box:conj-comm-3}
      \end{boxNested}
      \begin{boxNested}[1]
        \boxLine{Dokažimo $\phi$:}
        \boxLine[1]{Zaradi \boxRef{box:conj-comm-1} velja $\phi$.} \label{box:conj-comm-4}
      \end{boxNested}
    \end{boxFrame}
  \end{boxedProof}
  \end{center}
  %
  Celoten dokaz je vpeljava implikacije.
  %
  V vrsticah~\boxRef{box:conj-comm-2}--\boxRef{box:conj-comm-4} vidimo vpeljavo konjunkcije,
  v vrsticah \boxRef{box:conj-comm-3} in \boxRef{box:conj-comm-4} pa uporabo konjunkcije.
\end{zgled}

V dokazih boste videli frazo \nls{zadostuje dokazati}:
%
\begin{center}
  \begin{boxedProof}
    \begin{boxFrame}
      \boxLine{Dokažimo $\psi$:}
      \boxLine[1]{Zadostuje dokazati $\phi$:}
      \boxLine[2]{\brac{dokaz $\phi$}}
    \end{boxFrame}
  \end{boxedProof}
\end{center}
%
To je pravzaprav uporaba implikacije $\phi \lthen \psi$, ki pa jo je pisec zamolčal, ker meni, da je očitna. Bolj podrobno bi zapisali:
%
\begin{center}
  \begin{boxedProof}
    \begin{boxFrame}
      \boxLine{Dokažimo $\psi$:}
      \boxLine[1]{Ker velja $\phi \lthen \psi$, zadostuje dokazati $\phi$:}
      \boxLine[2]{\brac{dokaz $\phi$}}
    \end{boxFrame}
  \end{boxedProof}
\end{center}
%
Še vedno od bralca pričakujemo, da bo vedel, zakaj velja $\phi \lthen \psi$.
%
Običajno je to splošno znana implikacija, ki je nima smisla vedno znova ponavljati.


\paragraph*{Disjunkcija}
\label{sec:disjunkcija-pravila}

Disjunkcijo $\phi \lor \psi$ dokažemo tako, da dokažemo enega od disjunktov, zato imamo dve pravili vpeljave:
%
\begin{center}
  \begin{boxedProof}
    \begin{boxFrame}
      \boxLine{Dokažimo $\phi \lor \psi$:}
      \boxLine[1]{Dokažimo levi disjunkt $\phi$:}
      \boxLine[2]{\brac{dokaz $\phi$}}
    \end{boxFrame}
  \end{boxedProof}%
  \hspace{2em}%
  \begin{boxedProof}
    \begin{boxFrame}
      \boxLine{Dokažimo $\phi \lor \psi$:}
      \boxLine[1]{Dokažimo desni disjunkt $\psi$:}
      \boxLine[2]{\brac{dokaz $\psi$}}
    \end{boxFrame}
  \end{boxedProof}
\end{center}
%
Če vemo, da velja $\phi \lor \psi$, ne vemo pa, kateri od disjunktov velja, obravnavamo dva primera:
%
\begin{center}
  \begin{boxedProof}[reserved]
    \begin{boxFrame}
      \boxLineTag{$k$}{velja $\phi \lor \psi$}
      \boxGap
      \boxLine{Dokažimo $\rho$:}
      \boxLine[1]{Uporabimo ($k$) in obravnavamo dva primera:}
      \begin{boxNested}[1]
        \boxLine{Predpostavimo $\phi$.}
        \boxLine{Dokažimo $\rho$:}
        \boxLine[1]{\brac{dokaz $\rho$}}
      \end{boxNested}
      \begin{boxNested}[1]
        \boxLine{Predpostavimo $\psi$.}
        \boxLine{Dokažimo $\rho$:}
        \boxLine[1]{\brac{dokaz $\rho$}}
      \end{boxNested}
    \end{boxFrame}
  \end{boxedProof}
\end{center}
%
V prvem primeru smo izpeljali cilj pod predpostavko, da velja levi disjunkt~$\phi$, v drugem pa pod predpostavko, da velja desni disjunkt~$\psi$.

\begin{zgled}
  Obravnavajmo nekoliko bolj zapleten dokaz:
  %
  \begin{center}
  \begin{boxedProof}[numbered]
    \begin{boxFrame}
      \boxLine{Dokažimo $((p \lthen r) \land (q \lthen r)) \lthen (p \lor q \lthen r)$:}
      \boxLine[1]{Predpostavimo $(p \lthen r) \land (q \lthen r)$.\label{box:disj-elim-1}}
      \boxLine[1]{Dokažimo $p \lor q \lthen r$:\label{box:disj-elim-7}}
      \boxLine[2]{Predpostavimo $p \lor q$.\label{box:disj-elim-2}}
      \boxLine[2]{Dokažimo $r$:}
      \boxLine[3]{Uporabimo~\boxRef{box:disj-elim-2} in obravnavamo dva primera:\label{box:disj-elim-8}}
      \begin{boxNested}[3]
        \boxLine{Predpostavimo $p$.\label{box:disj-elim-3}}
        \boxLine{Dokažimo $r$:}
        \boxLine[1]{Zaradi~\boxRef{box:disj-elim-1} velja $p \lthen r$.\label{box:disj-elim-4}}
        \boxLine[1]{Zaradi~\boxRef{box:disj-elim-4} in~\boxRef{box:disj-elim-3} velja $r$.\label{box:disj-elim-9}}
      \end{boxNested}
      \begin{boxNested}[3]
        \boxLine{Predpostavimo $q$.\label{box:disj-elim-5}}
        \boxLine{Dokažimo $r$:}
        \boxLine[1]{Zaradi~\boxRef{box:disj-elim-1} velja $q \lthen r$.\label{box:disj-elim-6}}
        \boxLine[1]{Zaradi~\boxRef{box:disj-elim-6} in~\boxRef{box:disj-elim-5} velja $r$.\label{box:disj-elim-10}}
      \end{boxNested}
    \end{boxFrame}
  \end{boxedProof}
  \end{center}
  %
  Celoten dokaz je vpeljava implikacije.
  %
  V vrsticah~\boxRef{box:disj-elim-7}--\boxRef{box:disj-elim-10} vidimo še eno vpeljavo implikacije,
  v vrsticah~\boxRef{box:disj-elim-8}--\boxRef{box:disj-elim-10} pa uporabo disjunkcije.
  %
  V vrsticah \boxRef{box:disj-elim-4} in \boxRef{box:disj-elim-6} vidimo uporabo konjunkcije,
  v vrsticah \boxRef{box:disj-elim-9} in \boxRef{box:disj-elim-10} pa uporabo implikacije.
\end{zgled}

\paragraph*{Resnica}
\label{sec:resnica-pravila}

Resnice $\top$ ni treba dokazovati, zapišemo ">očitno"<:
%
\begin{center}
  \begin{boxedProof}
    \begin{boxFrame}
      \boxLine{Dokažimo $\top$:}
      \boxLine[1]{Očitno.}
    \end{boxFrame}
  \end{boxedProof}
\end{center}
%
Običajno $\top$ nastopi kot cilj, ko prvotni cilj poenostavimo v ekvivalentno obliko in dobimo~$\top$.
%
Na primer, cilj $12^2 + 12^2 < 17^2$ lahko z računanjem prevedemo v ekvivalentno izjavo $288 < 289$, ki je ekvivalentna~$\top$. S tem je dokaz zaključen, saj $\top$ ni treba dokazati.

Pravila uporabe za resnico ni.

\paragraph*{Neresnica}
\label{sec:neresnica-pravila}

Neresnica~$\bot$ nima pravila vpeljave. (To ne pomeni, da $\bot$ ni možno dokazati, saj lahko sledi iz protislovnih predpostavk.)
%
Pravilo uporabe za $\bot$ pravi, da iz neresnice sledi karkoli:
%
\begin{center}
  \begin{boxedProof}[reserved]
    \begin{boxFrame}
      \boxLineTag{$k$}{velja $\bot$}
      \boxGap
      \boxLine{zaradi~($k$) velja $\rho$}
    \end{boxFrame}
  \end{boxedProof}
\end{center}

\paragraph*{Ekvivalenca}
\label{sec:ekvivalenca-pravila}

Ekvivalenco $\phi \liff \psi$ dokažemo tako, da posebej dokažemo obe implikaciji:
%
\begin{center}
  \begin{boxedProof}
    \begin{boxFrame}
      \boxLine{Dokažimo $\phi \liff \psi$:}
      \begin{boxNested}
        \boxLine{Dokažimo $\phi \lthen \psi$:}
        \boxLine[1]{\brac{dokaz $\phi \lthen \psi$}}
      \end{boxNested}
      \begin{boxNested}
        \boxLine{Dokažimo $\psi \lthen \phi$:}
        \boxLine[1]{\brac{dokaz $\psi \lthen \phi$}}
      \end{boxNested}
    \end{boxFrame}
  \end{boxedProof}
\end{center}
%
Ekvivalenco uporabimo tako, da izpeljemo eno od implikacij:
%
%
\begin{center}
\begin{boxedProof}[reserved]
\begin{boxFrame}
  \boxLineTag{$k$}{velja $\phi \liff \psi$}
  \boxGap
  \boxLine{zaradi~($k$) velja $\phi \lthen \psi$}
\end{boxFrame}
\qquad\qquad
\end{boxedProof}
%
\begin{boxedProof}[reserved]
\begin{boxFrame}
  \boxLineTag{$k$}{velja $\phi \liff \psi$}
  \boxGap
  \boxLine{zaradi~($k$) velja $\psi \lthen \phi$}
\end{boxFrame}
\end{boxedProof}
\end{center}



\paragraph*{Negacija}
\label{sec:negacija-pravila}

Negacijo $\lnot \phi$ dokažemo tako, da predpostavimo~$\phi$ in dokažemo~$\bot$:
%
\begin{center}
  \begin{boxedProof}
    \begin{boxFrame}
      \boxLine{Dokažimo $\lnot \phi$:}
      \boxLine[1]{Predpostavimo $\phi$.}
      \boxLine[1]{Dokažimo~$\bot$:}
      \boxLine[2]{\brac{dokaz $\bot$}}
    \end{boxFrame}
  \end{boxedProof}
\end{center}
%
Običajno ne rečemo, da ">dokazujemo neresnico $\bot$"<, ampak da ">iščemo protislovje"<.

Pravilo uporabe za negacijo pravi, da iz $\phi$ in $\neg \phi$ sledi neresnica:
%
\begin{center}
  \begin{boxedProof}[reserved]
    \begin{boxFrame}
      \boxLineTag{$k$}{velja $\phi$}
      \boxGap
      \boxLineTag{$m$}{velja $\neg \phi$}
      \boxGap
      \boxLine{zaradi~($k$) in~($m$) velja $\bot$}
    \end{boxFrame}
  \end{boxedProof}
\end{center}
%
V dokazu običajno ne rečemo, da smo ">dokazali neresnico"<, ampak da smo ">izpeljali protislovje"<, se pravi izjavo in njeno negacijo.

\begin{zgled}
  Dokažimo, da iz protislovja sledi karkoli. Natančneje, iz $p$ in $\neg p$ sledi poljubna izjava~$q$:
  %
  \begin{center}
  \begin{boxedProof}[numbered]
    \begin{boxFrame}
      \boxLine{Dokažimo $p \land \neg p \lthen q$:}
      \boxLine[1]{Predpostavimo $p \land \neg p$.\label{box:ex-falso-1}}
      \boxLine[1]{Dokažimo $q$:}
      \boxLine[2]{Zaradi~\boxRef{box:ex-falso-1} velja $p$.\label{box:ex-falso-2}}
      \boxLine[2]{Zaradi~\boxRef{box:ex-falso-1} velja $\neg p$.\label{box:ex-falso-3}}
      \boxLine[2]{Zaradi~\boxRef{box:ex-falso-2} in~\boxRef{box:ex-falso-3} velja $\bot$.\label{box:ex-falso-4}}
      \boxLine[2]{Zaradi~\boxRef{box:ex-falso-4} velja $q$.\label{box:ex-falso-5}}
    \end{boxFrame}
  \end{boxedProof}
  \end{center}
  %
  Celoten dokaz je vpeljava implikacije.
  %
  V vrsticah~\boxRef{box:ex-falso-2} in~\boxRef{box:ex-falso-3} smo uporabili konjunkcijo,
  v vrstici~\boxRef{box:ex-falso-4} negacijo, v vrstici~\boxRef{box:ex-falso-5} pa neresnico.
\end{zgled}


\begin{zgled}
  \label{zgled:p-dvojna-negacija}
  Dokažimo, da iz $p$ sledi $\neg\neg p$:
  %
  \begin{center}
  \begin{boxedProof}[numbered]
    \begin{boxFrame}
      \boxLine{Dokažimo $p \lthen \neg\neg p$:}
      \boxLine[1]{Predpostavimo $p$.\label{box:dvojna-1}}
      \boxLine[1]{Dokažimo $\neg\neg p$:\label{box:dvojna-2}}
      \boxLine[2]{Predpostavimo $\neg p$.\label{box:dvojna-3}}
      \boxLine[2]{Dokažimo $\bot$:}
      \boxLine[3]{Zaradi~\boxRef{box:dvojna-1} in~\boxRef{box:dvojna-3} velja $\bot$.\label{box:dvojna-4}}
    \end{boxFrame}
  \end{boxedProof}
  \end{center}
  %
  Celoten dokaz je vpeljava implikacije.
  %
  V vrsticah~\boxRef{box:dvojna-2}--\boxRef{box:dvojna-4} vidimo vpeljavo negacije,
  v vrstici~\boxRef{box:dvojna-4} pa uporabo negacije.
\end{zgled}



\paragraph*{Univerzalni kvantifikator}
\label{sec:univerzalni-kvantifikator-pravila}

Izjavo $\all{x \in A} \phi(x)$ dokažemo tako, da vpeljemo spremenljivko $x \in A$ in dokažemo $\phi(x)$:
%
\begin{center}
  \begin{boxedProof}
    \begin{boxFrame}
      \boxLine{Dokažimo $\all{x \in A} \phi(x)$:}
      \boxLine[1]{Naj bo $x \in A$.}
      \boxLine[1]{Dokažimo $\phi(x)$:}
      \boxLine[2]{\brac{dokaz $\phi(x)$}}
    \end{boxFrame}
  \end{boxedProof}
\end{center}
%
\emph{Pozor:} spremenljivka $x$ mora biti \emph{sveža}, se pravi, da na tem mestu v dokazu ni veljavna. Če je $x$ že veljavna, izberemo svežo spremenljivko~$y$ in izjavo $\all{x \in A} \phi(x)$ najprej zapišemo kot $\all{y \in A} \phi(y)$.

Če vemo, da velja $\all{x \in A} \phi(x)$, in vemo, da je $\langle\text{izraz}\rangle \in A$, lahko izpeljemo $\phi(\langle\text{izraz}\rangle)$:
%
\begin{center}
  \begin{boxedProof}[reserved]
    \begin{boxFrame}
      \boxLineTag{$k$}{velja $\all{x \in A} \phi(x)$}
      \boxGap
      \boxLineTag{$m$}{velja $\langle\text{izraz}\rangle \in A$}
      \boxGap
      \boxLine{zaradi~($k$) in~($m$) velja $\phi(\langle\text{izraz}\rangle)$}
    \end{boxFrame}
  \end{boxedProof}
\end{center}

\begin{zgled}
  Sedaj lahko dokažemo, da izjava $\all{x \in \emptyset} \phi(x)$ velja, upoštevaje dejstvo, da za vsak~$x$ velja
  $\neg (x \in \emptyset)$:
  %
  \begin{center}
  \begin{boxedProof}[numbered]
    \begin{boxFrame}
      \boxLine{Dokažimo $\all{x \in \emptyset} \phi(x)$:}
      \boxLine[1]{Naj bo $x \in \emptyset$.\label{box:prazna-1}}
      \boxLine[1]{Dokažimo $\phi(x)$:}
      \boxLine[2]{Velja $\neg (x \in \emptyset)$.\label{box:prazna-2}}
      \boxLine[2]{Zaradi~\boxRef{box:prazna-1} in~\boxRef{box:prazna-2} velja $\bot$.\label{box:prazna-3}}
      \boxLine[2]{Zaradi~\boxRef{box:prazna-3} velja $\phi(x)$.\label{box:prazna-4}}
    \end{boxFrame}
  \end{boxedProof}
  \end{center}
  %
  Celoten dokaz je vpeljava univerzalnega kvantifikatorja.
  %
  V vrstici~\boxRef{box:prazna-3} smo uporabili negacijo, v vrstici~\boxRef{box:prazna-4} pa neresnico.
\end{zgled}


\paragraph*{Eksistenčni kvantifikator}
\label{sec:eksistencni-kvantifikator-pravila}

Izjavo $\some{x \in A} \phi(x)$ dokažemo tako, da podamo element množice~$A$, ki zadošča~$\phi$:
%
\begin{center}
  \begin{boxedProof}
    \begin{boxFrame}
      \boxLine{Dokažimo $\some{x \in A} \phi(x)$:}
      \boxLine[1]{Podamo $x \defeq \langle\text{izraz}\rangle$.}
      \begin{boxNested}[1]
        \boxLine{Dokažimo $\langle\text{izraz}\rangle \in A$:}
        \boxLine[1]{\brac{dokaz $\langle\text{izraz}\rangle \in A$}}
      \end{boxNested}
      \begin{boxNested}[1]
        \boxLine{Dokažimo $\phi(\langle\text{izraz}\rangle)$:}
        \boxLine[1]{\brac{dokaz $\phi(\langle\text{izraz}\rangle)$}}
      \end{boxNested}
    \end{boxFrame}
  \end{boxedProof}
\end{center}
%
Opomba: $\langle\text{izraz}\rangle$ sme vsebovati vse proste spremenljivke, ki so trenutno na voljo ($x$ \emph{ni} na
voljo).

Če vemo, da velja $\some{x \in A} \phi(x)$, smemo vpeljati spremenljivko $x \in A$, za katero velja $\phi(x)$:
%
\begin{center}
  \begin{boxedProof}[reserved]
    \begin{boxFrame}
      \boxLineTag{$k$}{velja $\some{x \in A} \phi(x)$}
      \boxGap
      \boxLine{Dokažimo $\rho$:}
      \boxLine[1]{Zaradi~($k$) imamo $x \in A$ in $\phi(x)$.}
      \boxLine[1]{\brac{dokaz $\rho$}}
    \end{boxFrame}
  \end{boxedProof}
\end{center}
%
\emph{Pozor:} spremenljivka $x$ mora biti sveža, se pravi, da na tem mestu v dokazu ni veljavna. Če je $x$ že veljavna, izberemo svežo spremenljivko~$y$ in izjavo $\some{x \in A} \phi(x)$ najprej zapišemo kot $\some{y \in A} \phi(y)$.

\begin{zgled}
  Dokažimo še, da velja $\neg \some{x \in \emptyset} \phi(x)$, upoštevaje dejstvo, da za vsak~$x$ velja
  $\neg (x \in \emptyset)$:
  %
  \begin{center}
  \begin{boxedProof}[numbered]
    \begin{boxFrame}
      \boxLine{Dokažimo $\neg \some{x \in \emptyset} \phi(x)$:}
      \boxLine[1]{Predpostavimo $\some{x \in \emptyset} \phi(x)$.\label{box:prazna-obstoj-1}}
      \boxLine[1]{Dokažimo $\bot$:} \boxLine[2]{Zaradi~\boxRef{box:prazna-obstoj-1} imamo $x \in \emptyset$ in $\phi(x)$.\label{box:prazna-obstoj-2}}
      \boxLine[2]{Zaradi \boxRef{box:prazna-obstoj-2} velja $x \in \emptyset$.\label{box:prazna-obstoj-5}}
      \boxLine[2]{Velja $\neg (x \in \emptyset)$.\label{box:prazna-obstoj-3}}
      \boxLine[2]{Zaradi~\boxRef{box:prazna-obstoj-5} in~\boxRef{box:prazna-obstoj-3} velja $\bot$.\label{box:prazna-obstoj-4}}
    \end{boxFrame}
  \end{boxedProof}
  \end{center}
  %
  Celoten dokaz je vpeljava negacije.
  %
  V vrstici~\boxRef{box:prazna-obstoj-2} smo uporabili eksistenčni kvantifikator, v
  vrstici~\boxRef{box:prazna-obstoj-5} konjunkcijo, v vrstici~\boxRef{box:prazna-obstoj-4} pa negacijo.
\end{zgled}


\section{Enakost in ekvivalenca}
\label{sec:enakost-ekvivalenca}

Pravila za enakost se glasijo:
%
\begin{itemize}
\item \df{refleksivnost}: $a = a$,
\item \df{simetričnost}: če $a = b$, potem $b = a$,
\item \df{tranzitivnost}: če $a = b$ in $b = c$, potem $a = c$,
\item \df{pravilo zamenjave}: če $a = b$ in $\phi(a)$, potem $\phi(b)$.
\end{itemize}
%
Zaenkrat enakosti ne bomo posvečali posebne pozornosti, saj jo v praksi študenti dobro obvladajo.

\begin{zgled}
  Pravilo za simetričnost enakosti sledi iz refleksivnosti in pravila zamenjave.
  Predpostavimo, da velja $a = b$, in obravnavajmo izjavo $\phi(x) \dfeq (x = a)$.
  Zaradi refleksivnosti velja $\phi(a)$, torej po pravilu zamenjave velja $\phi(b)$,
  kar je enako $b = a$.
\end{zgled}

\begin{vaja}
  Ali tudi tranzitivnost enakosti sledi iz refleksivnosti in pravila zamenjave?
\end{vaja}

Pravila za logično ekvivalenco so podobna pravilom za enakost:
%
\begin{itemize}
\item \df{refleksivnost}: $\phi \liff \phi$,
\item \df{simetričnost}: če $\phi \liff \psi$, potem $\psi \liff \phi$,
\item \df{tranzitivnost}: če $\phi \liff \psi$ in $\psi \liff \xi$, potem $\phi \liff \xi$,
\item \df{pravilo zamenjave}: če $\phi \liff \psi$, smemo $\phi$ nadomestiti s~$\psi$.
\end{itemize}
%
Zadnje pravilo dovoljuje, da lahko izjave poenostavljamo in predelujemo, če le pazimo, da ne kršimo logične ekvivalence.

\begin{vaja}
  S pravili sklepanja za ekvivalenco dokažite refleksivnost, simetričnost in tranzitivnost ekvivalence.
  Kako bi dokazali pravilo zamenjave? Upoštevajte, da smemo zamenjati izjavo z njej ekvivalentno kjerkoli, tudi znotraj druge izjave.
\end{vaja}

\paragraph*{Reševanje enačb}

V osnovni in srednji šoli nas učijo, da smemo na obeh straneh enačbe prišteti ali odšteti poljuben izraz, pomnožiti ali deliti smemo s poljubnim neničelnim izrazom, ipd. Od kod
izhajajo ta pravila? In kaj pomeni, da smo enačbo ">rešili"<? Rešitev enačbe
%
\begin{equation*}
  x^2 - 5 x + 6 = 0
\end{equation*}
%
običajno zapišemo takole:
%
\begin{equation*}
  x_1 = 2, \qquad x_2 = 3.
\end{equation*}
%
Kako naj to razumemo s stališča matematične logike? Treba je
pojasniti dvoje: kaj pomenita $x_1$ in $x_2$, saj v prvotni enačbi
nastopa spremenljivka $x$, ter kako naj razumemo vejico med izjavama
$x_1 = 2$ in $x_2 = 3$. Z indeksoma~$1$ in~$2$ štejemo rešitve enačbe
in sta v resnici nepotrebna,\footnote{Kako pa bi zapisali rešitve
  enačbe $x_1^2 - 5 x_1 + 6 = 0$?} na kar kaže tudi dejstvo, da
pišemo $x = \ldots$ in ne $x_1 = \ldots$, kadar je rešitev ena sama.
Torej bi lahko rešitev zapisali kot
%
\begin{equation*}
  x = 2, \quad x = 3.
\end{equation*}
%
Sedaj pa je tudi jasno, da bi namesto vejice morala stati disjunkcija,
se pravi
%
\begin{equation*}
  x = 2 \lor x = 3.
\end{equation*}
%
Začetna enačba in tako zapisana rešitev sta logično ekvivalentni:
%
\begin{equation*}
  x^2 - 5 x + 6 = 0 \iff
  x = 2 \lor x = 3.
\end{equation*}
%
Povzemimo: reševanje enačbe je postopek, s katerim dano enačbo $f(x) =
g(x)$ prevedemo v njeno \emph{logično ekvivalentno} obliko $x = a_1
\lor x = a_2 \lor \cdots \lor x = a_n$, iz katere so neposredno razvidne
rešitve enačbe.

Pravila za reševanje enačb torej niso nič drugega kot recepti, s
pomočjo katerih enačbo predelamo v njeno \emph{ekvivalentno} obliko, ki
je korak bližje končni obliki, v kateri bi radi zapisali rešitev. To
pojasnjuje srednješolska pravila za reševanje enačb. Na primer, za
realna števila $a, b, c \in \RR$ vedno velja
%
\begin{equation*}
  a = b \lthen c \cdot a = c \cdot b,
\end{equation*}
%
medtem ko obratna implikacija
%
\begin{equation*}
  c \cdot a = c \cdot b \lthen a = b
\end{equation*}
%
za splošne $a$ in $b$ velja le v primeru, ko je $c \neq 0$. Ker pri
reševanju enačb potrebujemo implikacijo v obe smeri, srednješolce
učimo, da smejo enačbo množiti samo z od nič različnimi števili.

\begin{vaja}
  Ali smemo obe strani enačbe kvadrirati?
\end{vaja}

\begin{vaja}
  Denimo, da v postopku reševanja enačbe pomotoma naredimo korak,
  ki ni ekvivalenca, ampak le implikacija:
  %
  \begin{align*}
    f_1(x) = g_1(x) &\liff \\
    \vdots & \\
    f_k(x) = g_k(x) &\lthen \tag{napačen korak}\\
    f_{k+1}(x) = g_{k+1}(x) &\liff \\
    \vdots & \\
    x = a_1 \lor \cdots \lor x = a_n&
  \end{align*}
  %
  Smo dobili preveč ali premalo rešitev prvotne enačbe?
\end{vaja}


\section{Pravilo o izključenem tretjem}
\label{sec:lem}

Pravilo o izključenem tretjem se glasi:
%
\begin{equation*}
  \nls{Za vsako izjavo $\phi$ velja $\phi \lor \neg\phi$.}
\end{equation*}
%
Se pravi, da je vsaka izjava bodisi resnična bodisi neresnična.
%
Torej ni ">tretje možnosti"< za resničnostno vrednost izjave $\phi$,
od koder izhaja ime pravila.

Pravilo o izključenem tretjem omogoča \emph{posredne} dokaze izjav, od
katerih je najbolj znano \df{dokazovanje s protislovjem}:
%
\begin{center}
  \begin{boxedProof}
    \begin{boxFrame}
      \boxLine{Dokažimo $\phi$ s protislovjem:}
      \boxLine[1]{Predpostavimo $\neg\phi$.}
      \boxLine[1]{Dokažimo~$\bot$:}
      \boxLine[2]{\brac{dokaz $\bot$}}
    \end{boxFrame}
  \end{boxedProof}
\end{center}
%
Dokaz s protislovjem sloni na dejstvu, da lahko dokažemo $\neg\neg\phi \lthen \phi$, vendar pri tem potrebujemo pravilo o izključenem tretjem:
%
\begin{center}
  \begin{boxedProof}[numbered]
    \begin{boxFrame}
      \boxLine{Dokažimo $\neg\neg\phi \lthen \phi$:}
      \boxLine[1]{Predpostavimo $\neg\neg\phi$.\label{box:lem-dvojna-1}}
      \boxLine[1]{Dokažimo $\phi$:}
      \boxLine[2]{Ker velja $\phi \lor \neg\phi$, obravnavamo dva primera:}
      \begin{boxNested}[2]
        \boxLine{Predpostavimo $\phi$.\label{box:lem-dvojna-2}}
        \boxLine{Dokažimo $\phi$:}
        \boxLine[1]{Zaradi~\boxRef{box:lem-dvojna-2} velja $\phi$.}
      \end{boxNested}
      \begin{boxNested}[2]
        \boxLine{Predpostavimo $\neg\phi$.\label{box:lem-dvojna-3}}
        \boxLine{Dokažimo $\phi$:}
        \boxLine[1]{Zaradi~\boxRef{box:lem-dvojna-3} in~\boxRef{box:lem-dvojna-1} velja $\bot$.\label{box:lem-dvojna-4}}
        \boxLine[1]{Zaradi~\boxRef{box:lem-dvojna-4} velja $\phi$.}
      \end{boxNested}
    \end{boxFrame}
  \end{boxedProof}
\end{center}
%
Velja tudi obrat: pravilo o izključenem tretjem lahko izpeljemo z dokazom s protislovjem:
%
\begin{center}
  \begin{boxedProof}[numbered]
    \begin{boxFrame}
      \boxLine{Dokažimo $\phi \lor \neg\phi$ s protislovjem:}
      \boxLine[1]{Predpostavimo $\neg (\phi \lor \neg\phi)$.\label{box:lem-tretja-1}}
      \boxLine[1]{Dokažimo $\bot$:}
      \boxLine[2]{Velja $\phi \lor \neg\phi$:\label{box:lem-tretja-5}}
      \begin{boxNested}[2]
        \boxLine{Dokažimo desni disjunkt $\neg\phi$:}
        \boxLine[1]{Predpostavimo $\phi$.\label{box:lem-tretja-2}}
        \boxLine[1]{Dokažimo $\bot$:}
        \boxLine[2]{Zaradi~\boxRef{box:lem-tretja-2} velja $\phi \lor \neg\phi$.\label{box:lem-tretja-3}}
        \boxLine[2]{Zaradi~\boxRef{box:lem-tretja-3} in~\boxRef{box:lem-tretja-1} velja $\bot$.}
      \end{boxNested}
      \boxLine[2]{Zaradi~\boxRef{box:lem-tretja-5} in~\boxRef{box:lem-tretja-1} velja $\bot$.}
    \end{boxFrame}
  \end{boxedProof}
\end{center}

Pravilo o izključenem tretjem omogoča še nekatere druge načine dokazovanja izjav.
Najprej, implikacijo $\phi \lthen \psi$ smemo zamenjati s \df{kontrapozicijo} $\neg \psi \lthen \neg \phi$, saj sta ekvivalentni.
%
V eno smer dokaz ekvivalence poteka takole:
%
\begin{center}
  \begin{boxedProof}[numbered]
    \begin{boxFrame}
      \boxLine{Dokažimo $(\phi \lthen \psi) \lthen (\neg\psi \lthen \neg\phi)$:}
      \boxLine[1]{Predpostavimo $\phi \lthen \psi$.\label{box:kontrapozicija-1}}
      \boxLine[1]{Dokažimo $\neg\psi \lthen \neg\phi$:}
      \boxLine[2]{Predpostavimo $\neg\psi$.\label{box:kontrapozicija-2}}
      \boxLine[2]{Dokažimo $\neg\phi$:}
      \boxLine[3]{Predpostavimo $\phi$.\label{box:kontrapozicija-3}}
      \boxLine[3]{Dokažimo $\bot$:}
      \boxLine[4]{Zaradi~\boxRef{box:kontrapozicija-1} in~\boxRef{box:kontrapozicija-3} velja $\psi$.\label{box:kontrapozicija-4}}
      \boxLine[4]{Zaradi~\boxRef{box:kontrapozicija-4} in~\boxRef{box:kontrapozicija-2} velja $\bot$.}
    \end{boxFrame}
  \end{boxedProof}
\end{center}
%
V drugo smer potrebujemo dokaz s protislovjem:
%
\begin{center}
  \begin{boxedProof}[numbered]
    \begin{boxFrame}
      \boxLine{Dokažimo $(\neg\psi \lthen \neg\phi) \lthen (\phi \lthen \psi)$:}
      \boxLine[1]{Predpostavimo $\neg\psi \lthen \neg\phi$.\label{box:kontrapozicija-5}}
      \boxLine[1]{Dokažimo $\phi \lthen \psi$:}
      \boxLine[2]{Predpostavimo $\phi$.\label{box:kontrapozicija-6}}
      \boxLine[2]{Dokažimo $\psi$ s protislovjem:}
      \boxLine[3]{Predpostavimo $\neg\psi$.\label{box:kontrapozicija-7}}
      \boxLine[3]{Dokažimo $\bot$:}
      \boxLine[4]{Zaradi~\boxRef{box:kontrapozicija-5} in~\boxRef{box:kontrapozicija-7} velja $\neg\phi$.\label{box:kontrapozicija-8}}
      \boxLine[4]{Zaradi~\boxRef{box:kontrapozicija-6} in~\boxRef{box:kontrapozicija-8} velja $\bot$.}
    \end{boxFrame}
  \end{boxedProof}
\end{center}
%
Dalje, disjunkcijo $\phi \lor \psi$ smemo zamenjati z implikacijo $\neg \phi \lthen \psi$.
%
Ekvivalenco $(\phi \lor \psi) \liff (\neg\phi \lthen \psi)$ v eno smer dokažemo takole:
%
\begin{center}
  \begin{boxedProof}[numbered]
    \begin{boxFrame}
      \boxLine{Dokažimo $(\phi \lor \psi) \lthen (\neg\phi \lthen \psi)$:}
      \boxLine[1]{Predpostavimo $\phi \lor \psi$.\label{box:disjunkcija-implikacija-1}}
      \boxLine[1]{Dokažimo $\neg\phi \lthen \psi$:}
      \boxLine[2]{Predpostavimo $\neg\phi$.\label{box:disjunkcija-implikacija-2}}
      \boxLine[2]{Dokažimo $\psi$:}
      \boxLine[3]{Uporabimo~\boxRef{box:disjunkcija-implikacija-1} in obravnavamo dva primera:}
      \begin{boxNested}[3]
        \boxLine{Predpostavimo $\phi$.\label{box:disjunkcija-implikacija-3}}
        \boxLine{Dokažimo $\psi$:}
        \boxLine[1]{Zaradi~\boxRef{box:disjunkcija-implikacija-3} in~\boxRef{box:disjunkcija-implikacija-2} velja $\bot$.\label{box:disjunkcija-implikacija-4}}
        \boxLine[1]{Zaradi~\boxRef{box:disjunkcija-implikacija-4} velja $\psi$.}
      \end{boxNested}
      \begin{boxNested}[3]
        \boxLine{Predpostavimo $\psi$.\label{box:disjunkcija-implikacija-5}}
        \boxLine{Dokažimo $\psi$:}
        \boxLine[1]{Zaradi~\boxRef{box:disjunkcija-implikacija-5} velja $\psi$.}
      \end{boxNested}
    \end{boxFrame}
  \end{boxedProof}
\end{center}
%
V drugo smer potrebujemo pravilo o izključenem tretjem:
%
\begin{center}
  \begin{boxedProof}[numbered]
    \begin{boxFrame}
      \boxLine{Dokažimo $(\neg\phi \lthen \psi) \lthen (\phi \lor \psi)$:}
      \boxLine[1]{Predpostavimo $\neg\phi \lthen \psi$.\label{box:disjunkcija-implikacija-6}}
      \boxLine[1]{Dokažimo $\phi \lor \psi$:}
      \boxLine[2]{Ker velja $\phi \lor \neg\phi$, obravnavamo dva primera:}
      \begin{boxNested}[2]
        \boxLine{Predpostavimo $\phi$.\label{box:disjunkcija-implikacija-7}}
        \boxLine{Dokažimo $\phi \lor \psi$:}
        \boxLine[1]{Dokažimo levi disjunkt $\phi$:}
        \boxLine[2]{Zaradi~\boxRef{box:disjunkcija-implikacija-7} velja $\phi$.}
      \end{boxNested}
      \begin{boxNested}[2]
        \boxLine{Predpostavimo $\neg\phi$.\label{box:disjunkcija-implikacija-8}}
        \boxLine{Dokažimo $\phi \lor \psi$:}
        \boxLine[1]{Dokažimo desni disjunkt $\psi$:}
        \boxLine[2]{Zaradi~\boxRef{box:disjunkcija-implikacija-6} in~\boxRef{box:disjunkcija-implikacija-8} velja $\psi$.}
      \end{boxNested}
    \end{boxFrame}
  \end{boxedProof}
\end{center}

Pravilo o izključenem tretjem je koristno tudi za poenostavljanje negacij
s pomočjo naslednjih ekvivalenc, ki povedo, kako~$\lnot$ ">porinemo"< v izjavo:
%
\begin{align*}
  \lnot (\phi \land \psi) &\iff \lnot\phi \lor \lnot\psi \\
  \lnot (\phi \lor \psi) &\iff \lnot\phi \land \lnot\psi \\
  \lnot (\phi \lthen \psi) &\iff \phi \land \lnot\psi \\
  \lnot (\lnot \phi) &\iff \phi \\
  \lnot (\all{x \in S} \phi(x)) &\iff \some{x \in S} \lnot\phi(x) \\
  \lnot (\some{x \in S} \phi(x)) &\iff \all{x \in S} \lnot\phi(x)
\end{align*}
%
Dokažimo le predzadnjo ekvivalenco, ostale pa prepustimo za vajo (zadnja ekvivalenca velja brez uporabe izključenega tretjega).
%
V eno smer dokaz poteka takole:
%
\begin{center}
  \begin{boxedProof}[numbered]
    \begin{boxFrame}
      \boxLine{Dokažimo $(\some{x \in S} \lnot\phi(x)) \lthen \lnot (\all{x \in S} \phi(x))$:}
      \boxLine[1]{Predpostavimo $\some{x \in S} \lnot\phi(x)$.\label{box:negacija-vsak-1}}
      \boxLine[1]{Dokažimo $\lnot (\all{x \in S} \phi(x))$:}
      \boxLine[2]{Predpostavimo $\all{x \in S} \phi(x)$.\label{box:negacija-vsak-2}}
      \boxLine[2]{Dokažimo $\bot$:}
      \boxLine[3]{Zaradi~\boxRef{box:negacija-vsak-1} imamo $x \in S$ in $\lnot\phi(x)$.\label{box:negacija-vsak-3}}
      \boxLine[3]{Zaradi~\boxRef{box:negacija-vsak-3} velja $x \in S$.\label{box:negacija-vsak-4}}
      \boxLine[3]{Zaradi~\boxRef{box:negacija-vsak-3} velja $\lnot\phi(x)$.\label{box:negacija-vsak-5}}
      \boxLine[3]{Zaradi~\boxRef{box:negacija-vsak-2} in~\boxRef{box:negacija-vsak-4} velja $\phi(x)$.\label{box:negacija-vsak-6}}
      \boxLine[3]{Zaradi~\boxRef{box:negacija-vsak-6} in~\boxRef{box:negacija-vsak-5} velja $\bot$.}
    \end{boxFrame}
  \end{boxedProof}
\end{center}
%
V drugo smer potrebujemo dokaz s protislovjem:
%
\begin{center}
  \begin{boxedProof}[numbered]
    \begin{boxFrame}
      \boxLine{Dokažimo $\lnot (\all{x \in S} \phi(x)) \lthen \some{x \in S} \lnot\phi(x)$:}
      \boxLine[1]{Predpostavimo $\lnot (\all{x \in S} \phi(x))$.\label{box:negacija-vsak-7}}
      \boxLine[1]{Dokažimo $\some{x \in S} \lnot\phi(x)$ s protislovjem:}
      \boxLine[2]{Predpostavimo $\lnot (\some{x \in S} \lnot\phi(x))$.\label{box:negacija-vsak-8}}
      \boxLine[2]{Dokažimo $\bot$:}
      \boxLine[3]{Velja $\all{x \in S} \phi(x)$:\label{box:negacija-vsak-9}}
      \begin{boxNested}[3]
        \boxLine{Naj bo $x \in S$.\label{box:negacija-vsak-10}}
        \boxLine{Dokažimo $\phi(x)$ s protislovjem:}
        \boxLine[1]{Predpostavimo $\lnot\phi(x)$.\label{box:negacija-vsak-11}}
        \boxLine[1]{Dokažimo $\bot$:}
        \boxLine[2]{Zaradi~\boxRef{box:negacija-vsak-10} in~\boxRef{box:negacija-vsak-11} velja $\some{x \in S} \lnot\phi(x)$.\label{box:negacija-vsak-12}}
        \boxLine[2]{Zaradi~\boxRef{box:negacija-vsak-12} in~\boxRef{box:negacija-vsak-8} velja $\bot$.}
      \end{boxNested}
      \boxLine[3]{Zaradi~\boxRef{box:negacija-vsak-9} in~\boxRef{box:negacija-vsak-7} velja $\bot$.}
    \end{boxFrame}
  \end{boxedProof}
\end{center}

Posrednih načinov dokazovanja, ki smo jih predstavili tu, se poslužujemo le, ko neposredni dokaz ni možen ali pa je zapleten, saj nam ne povedo prav veliko o dokazani izjavi. Na primer, v neposrednem dokazu izjave $\some{x \in A} \phi(x)$ podamo element~$a \in A$, ki zadošča $\phi(a)$, s posrednim dokazom ekvivalentne izjave $\neg \all{x \in A} \neg\phi(x)$ pa samo dokažemo, da ni možno, da tak element ne bi obstajal.

\begin{zgled}
  Pokažimo, kako zgoraj navedene ekvivalence uporabimo, ko dokazujemo negacijo izjave.
  Denimo, da bi radi ovrgli izjavo
  % 
  \begin{quote}
    \nls{Vsako zaporedje pozitivnih realnih števil ima limito~$0$.}
  \end{quote}
  %
  Najprej jo zapišemo v simbolni obliki in pred njo postavimo negacijo, saj izjavo ovržemo tako, da dokažemo njeno negacijo:
  %
  \begin{equation}
    \label{eq:pozitivno-limita-0}
    \neg \all{a \in \RR^{\NN}}{\text{$(a_n)_n$ pozitivno zaporedje}
      \lthen \text{$0$ je limita zaporedja $(a_n)_n$}}.
  \end{equation}
  % 
  Zgornja pravila za računanje negacije nam povedo, da se
  $\lnot\forall$ spremeni v $\exists\lnot$ in da se nato implikacija
  oblike $\phi \lthen \psi$ spremeni v $\phi \land \lnot\psi$. Tako
  dobimo ekvivalentno obliko~\eqref{eq:pozitivno-limita-0}:
  % 
  \begin{equation*}
    \some{a \in \RR^{\NN}}{\text{$(a_n)_n$ pozitivno zaporedje}
      \land \lnot (\text{$0$ je limita zaporedja $(a_n)_n$})}.
  \end{equation*}
  % 
  To preberemo z besedami:
  % 
  \begin{quote}
    \nls{Obstaja tako zaporedje $(a_n)_n$, da je $(a_n)_n$ zaporedje
    pozitivnih števil in da $0$ ni limita zaporedja $(a_n)_n$.}
  \end{quote}
  %
  Če se še malo potrudimo, preberemo bolj razumljivo:
  % 
  \begin{quote}
    \nls{Obstaja tako zaporedje pozitivnih realnih števil, da $0$ ni
    njegova limita.}
  \end{quote}
  %
  S tem še nismo končali, saj je \nls{Število $0$ ni limita
  zaporedja $(a_n)_n$} negacija izjave \nls{$0$ je limita zaporedja
  $(a_n)_n$}, kar v simbolni obliki zapišemo
  % 
  \begin{equation}
    \label{eq:limita-0}
    \neg \all{\epsilon > 0}
      \some{m \in \NN}
        \all{n \geq m}
          |a_n - 0| < \epsilon.
  \end{equation}
  % 
  Z zgornjimi pravili za negiranje operacijo $\lnot$ postopoma ">porivamo"<
  navznoter:
  % 
  % 
  \begin{align*}
    \lnot \all{\epsilon > 0} \some{m \in \NN} \all{n \geq m} |a_n
          - 0| < \epsilon & \iff
    \\
    \some{\epsilon > 0} \lnot \some{m \in \NN} \all{n \geq m}
          |a_n - 0| < \epsilon &\iff
    \\
    \some{\epsilon > 0} \all{m \in \NN} \lnot \all{n \geq m}
          |a_n - 0| < \epsilon &\iff
    \\
    \some{\epsilon > 0} \all{m \in \NN} \some{n \geq m}
          \lnot (|a_n - 0| < \epsilon) &\iff
    \\
    \some{\epsilon > 0} \all{m \in \NN} \some{n \geq m}
          |a_n - 0| \geq \epsilon &\iff
    \\
    \some{\epsilon > 0} \all{m \in \NN} \some{n \geq m}
          |a_n| \geq \epsilon &\iff
    \\
    \some{\epsilon > 0} \all{m \in \NN} \some{n \geq m}
          a_n \geq \epsilon.
  \end{align*}
  %
  Ugotovili smo, da iščemo zaporedje $(a_n)_n$ pozitivnih števil, ki zadošča
  %
  \begin{quote}
    \nls{Obstaja tak $\epsilon > 0$, da za vsak $m \in \NN$ obstaja $n \geq m$,
      da je $a_n \geq \epsilon$.}
  \end{quote}
  %
  Kako najdemo zaporedje, ki zadošča temu pogoju?
  %
  Poskusimo kar naivno: vstavimo neko konkretno vrednost za $\epsilon$, na primer~$42$:
  %
  \begin{quote}
    \nls{Za vsak $m \in \NN$ obstaja $n \geq m$, da je $a_n \geq 42$.}
  \end{quote}
  %
  Potrebujemo zaporedje z neskončno mnogo členi, ki so vsaj~$42$, na primer $a_n \defeq 43$.
  %
  Res, če je $m \in \NN$ poljuben, lahko vzamemo kar $n = m$, saj velja $a_n = a_m = 43 > 42 = \epsilon$.

  Povzemimo odgovor: ni res, da ima vsako zaporedje pozitivnih realnih števil limito~$0$, saj zaporedje $a_n \defeq 43$ nima limite~$0$.
\end{zgled}

%%%%%%%%%%%%%%%%%%%%%%%%%%%%%%%%%%%%%%%%%%%%%%%%%%%%%%%%%%%%%%%%%%%%%%

\section{Brez izgube za splošnost}

Poglavje zaključimo še s tehniko dokazovanja, ki jo boste srečali v praksi in ki od bralca zahteva malo več znanja kot osnovna pravila sklepanja.

Fraza \nls{brez izgube za splošnost \dots} nakazuje, da dokaz obravnava le poseben primer, iz katerega sledi cilj, oziroma je preostanek dokaza zelo podoben posebnemu primeru. Bralec mora sam razbrati, zakaj je tako.
Podajmo primer.

\begin{trditev}
  \label{izrek:abc-vsota-razlik-soda}
  Za vsa cela števila $a$, $b$ in $c$ je $|a-b|+|b-c|+|c-a|$ sodo število.
\end{trditev}

\begin{proof}
  Brez izgube za splošnost smemo predpostaviti $a \geq b \geq c$.
  Tedaj velja
  %
  \begin{equation*}
    |a-b| + |b-c| + |c-a| = (a - b) + (b - c) - (c - a) = 2 (a - c),
  \end{equation*}
  %
  kar je sodo število.
\end{proof}

Zakaj je avtor dokaza predpostavil, da velja $a \geq b \geq c$? 
%
S frazo \nls{brez izgube za splošnost} nam sporoča, da ni treba obravnavati vseh šestih primerov
%
\begin{xalignat*}{3}
  & a \geq b \geq c,&
  & a \geq c \geq b,&
  & b \geq a \geq c,\\
  & b \geq c \geq a,&
  & c \geq a \geq b,&
  & c \geq b \geq a,
\end{xalignat*}
%
ker lahko vsakega od njih prevedemo na $a \geq b \geq c$.
%
In res, ker števila~$a$, $b$ in~$c$ v izrazu $|a-b|+|b-c|+|c-a|$ nastopajo \emph{simetrično}, se vrednost izraza ne spremeni, če jih premešamo.
%
Denimo, ko zamenjamo $a$ in $b$, dobimo $|b-a|+|a-c|+|c-b|$, kar je
enako prvotnemu izrazu $|a-b|+|b-c|+|c-a|$.
%
Števila torej uredimo po padajočem vrstnem redu~$a \geq b \geq c$ in obravnavamo le en primer.


%%% Local Variables:
%%% mode: latex
%%% TeX-master: "lmn"
%%% End:
